% BANKS-SOURCE: pdf=39 print=29 kind=implicit-exercise id=I-CH03-014
\begin{exercise}[Why the Legendre transform makes trees]{I-CH03-014}
Let
\[
 W[J]=\int\dd^D x\,\phi(x)J(x)+\Gamma[\phi],
 \qquad \phi=\frac{\delta W}{\delta J},
 \qquad J=-\frac{\delta\Gamma}{\delta\phi}.
\]
Treat all second derivatives as integral kernels.  Derive
\(W_2=-\Gamma_2^{-1}\), differentiate it once, and obtain
\[
 W_3(x,y,z)=\int\dd^D a\,\dd^D b\,\dd^D c\,
 W_2(x,a)W_2(y,b)W_2(z,c)\Gamma_3(a,b,c).
\]
Differentiate again and display the
four-point answer as one $\Gamma_4$ tree and the three exchange trees made
from two $\Gamma_3$ vertices.  Prove by induction that every connected
$W_n$ is a sum of trees with full propagators $W_2$ on the branches and 1PI
vertices $\Gamma_k$.  Explain why the differentiation step generates every
tree with a newly labeled external leg exactly once.
\end{exercise}
